\subsection*{主机、Tileflow 与设备代码}

普通 C++ 代码还可以包含设备辅助函数。因此，完整示例中常见三个区域：设备辅助函数、调用它的 \co@__co__@ 函数，以及主机侧的 \co@main()@。下面按这个顺序列出它们。

\begin{lstlisting}[language=croqtile,caption={单份源文件中的三个区域}]
__device__ void device_function(...) {
  // High-performance device kernel implementation
}
// Tileflow program: orchestrating data movement
__co__ void croqtile_function(...) {
  // ... croqtile code ...
  device_function(...);
  // ...
}
// Host program: typically runs on CPU
void main() {
  // ... prepare data ...
  croqtile_function(...);
  // ...
}
\end{lstlisting}

程序从主机侧开始执行。主机代码使用普通 C++，负责准备输入、调用 Croqtile 函数并处理结果。

\co@__co__@ 函数构成 Tileflow 程序，负责组织并行工作和数据搬运。它将数据放到合适的缓冲区，安排设备端计算，并把结果交回调用方。

\co@__device__@ 辅助函数执行设备端计算。Croqtile 编译器主要翻译 \co@__co__@ 区域，普通 C++ 的主机代码和设备代码交给后续工具链处理。\textless{}br\textgreater{}\textless{}br\textgreater{}\textbf{提示：} \co@__device__@ 主要用于在必要时复用特定硬件的函数，简化开发。对于 Croqtile 已提供较完善支持的硬件，通常直接使用语言提供的计算与数据搬运能力即可，无需通过设备辅助函数自行补充底层实现。
